\documentclass[12pt, a4paper]{article}

% Including essential packages for formatting and structure
\usepackage[utf8]{inputenc}
\usepackage[T1]{fontenc}
\usepackage{geometry}
\usepackage{graphicx}
\usepackage{booktabs}
\usepackage{amsmath}
\usepackage{hyperref}
\usepackage{url}
\usepackage{fancyhdr}
\usepackage{titling}
\usepackage[vietnamese]{babel}

% Configure URL line breaks
\Urlmuskip=0mu plus 1mu

% Configuring page geometry for proper margins
\geometry{margin=1in}

% Setting up headers and footers
\pagestyle{fancy}
\fancyhf{}

\title{Classification of Audio Embeddings}
\author{Turkey Detection Team}
\date{June 01, 2025}

\fancyhead[L]{\leftmark}
\fancyhead[R]{\thepage}
\fancyfoot[C]{\thetitle}

% Configuring title and author for the cover page
\pretitle{\begin{center}\LARGE\bfseries}
\posttitle{\par\end{center}\vskip 0.5em}
\preauthor{\begin{center}\large}
\postauthor{\end{center}}

% Setting the font to Latin Modern for consistency
\usepackage{lmodern}

\begin{document}

% Creating the cover page
\begin{titlepage}
    \centering
    \vspace*{2cm}

    \vspace{1.5cm}

    \vspace{1.5cm}
    
    % Report title
    \Huge \textbf{Classification of Audio Embeddings} \\
    
    \vspace{1cm}
    
    % Subject
    \Large \textbf{Subject: Introduction to Machine Learning} \\
        
    \vspace{0.5cm}

    \includegraphics[scale=.40]{img/hcmus-logo.png}\\[1cm]
    
    \vspace{0.5cm}
    
    % Supervised teachers
    \large \textbf{Supervised by: Bùi Duy Đăng, Nguyễn Thanh Tình} \\

    \vspace{0.5cm}

    % Location
    \large \textbf{Location: Thành phố Hồ Chí Minh, Việt Nam} \\

    \vspace{0.5cm}
    
    % School name
    \large \textbf{Trường Đại học Khoa Học Tự Nhiên, tp Hồ Chí Minh} \\
    
    \vspace{1cm}
    
    % Group information
    \large \textbf{Prepared by: Group 30} \\

    % Date
    \small \text{Date: June 01, 2025} \\
    
    \vfill
    
\end{titlepage}

\newpage

% Group Information Section
\section{Group Information}
This project was completed by the group \textbf{Group 30}. The team members are listed below:

\begin{table}[h]
    \centering
    \begin{tabular}{ll}
        \toprule
        \textbf{Name} & \textbf{Student ID} \\
        \midrule
        Bùi Kim Phúc & 21120112 \\
        Lê Hoàng Sơn & 21120127 \\
        \bottomrule
    \end{tabular}
    \caption{Group Members}
    \label{tab:group_members}
\end{table}

% About the Project Section
\section{About the Project}
The objective of this project is to develop a machine learning model to classify audio clips based on the presence of turkey sounds. The input data consists of audio embeddings, each represented as a matrix of shape $[10, 128]$, where 10 is the number of frames and 128 is the dimensionality of each frame’s feature vector. These embeddings are extracted from audio clips and provided in a JSON file (\texttt{train.json}). The task is a binary classification problem, where the model predicts whether an audio clip contains turkey sounds (\texttt{is\_turkey = 1}) or not (\texttt{is\_turkey = 0}). The classification is performed using a MLP model using Pytorch framework, leveraging the flattened embeddings (size 1280) as input features.
Finally, the model's performance is evaluated on a test set, and the results are submitted in a CSV file (\texttt{submission.csv}) for Kaggle competition.
% Data Preprocessing Section
\section{Data Preprocessing}
The dataset was sourced from a JSON file (\texttt{train.json}) containing audio embeddings and binary labels. The preprocessing steps included:

\begin{itemize}
    \item Loading the JSON data into a pandas DataFrame.
    \item Extracting \texttt{audio\_embedding} (shape [10, 128]) and \texttt{is\_turkey} (binary labels: 0 or 1).
    \item Flattening each audio embedding to a 1D array of size 1280 (10 $\times$ 128).
    \item Padding or truncating embeddings to ensure a consistent size.
    \item The data contains a total of 1195 samples and is split into 836 training samples (70\%), 179 validation samples (15\%), and 180 test samples (15\%) with shuffling to ensure randomness.
\end{itemize}

% Model Description Section
\section{Model Description}
This project was conducted in three phases, employing different model architectures:

\subsection{Phase 1: Traditional Machine Learning Models}
Initially, we implemented two classical machine learning models using Scikit-learn:

\begin{itemize}
    \item \textbf{Logistic Regression}:
        \begin{itemize}
            \item \textbf{Algorithm}: Logistic regression with the \texttt{liblinear} solver, suitable for small datasets.
            \item \textbf{Input Features}: Flattened audio embeddings of size 1280 (10 frames $\times$ 128 dimensions).
            \item \textbf{Output}: Binary classification (0 or 1) for the \texttt{is\_turkey} label.
            \item \textbf{Hyperparameters}: C=1.0, max\_iter=1000, random\_state=42.
        \end{itemize}
    \item \textbf{Random Forest}:
        \begin{itemize}
            \item \textbf{Algorithm}: Random forest classifier, an ensemble method using multiple decision trees.
            \item \textbf{Input Features}: Flattened audio embeddings of size 1280 (10 frames $\times$ 128 dimensions).
            \item \textbf{Output}: Binary classification (0 or 1) for the \texttt{is\_turkey} label.
            \item \textbf{Hyperparameters}: n\_estimators=100, max\_depth=None, random\_state=42.
        \end{itemize}
\end{itemize}

\subsection{Phase 2: Neural Network Approach}
In the second phase, we implemented a more sophisticated neural network model using PyTorch:

\begin{itemize}
    \item \textbf{Multilayer Perceptron (MLP)}:
        \begin{itemize}
            \item \textbf{Architecture}: A flexible neural network with configurable hidden layers.
            \item \textbf{Input Layer}: Flattens the input embeddings from [10, 128] to a vector of size 1280.
            \item \textbf{Hidden Layers}: Multiple fully connected layers with ReLU activation functions.
            \item \textbf{Regularization}: Dropout layers with 0.3 dropout rate to prevent overfitting.
            \item \textbf{Output Layer}: A fully connected layer producing logits for binary classification.
        \end{itemize}
\end{itemize}

MLP model compared to models in Phase 1 offers several advantages:
\begin{itemize}
    \item More complex pattern recognition through non-linear transformations using Rectified Linear Units (ReLU) activation functions.
    \item Better handling of the high-dimensional audio embedding features
    \item Enhanced regularization strategies to prevent overfitting
\end{itemize}

% Platform Section
\section{Platform}
The project was developed and executed on the following platform:

\begin{itemize}
    \item \textbf{Google Colab}: The primary development environment, providing a cloud-based Jupyter notebook interface with GPU support for efficient model training and evaluation.
    \item \textbf{Google Drive}: Used for storing and accessing the dataset (\texttt{train.json}) and saving the submission file (\texttt{submission.csv}).
    \item \textbf{Libraries}: Pytorch framework for model implementation, pandas for data processing, NumPy for numerical operations, and other standard Python libraries.
    \item \textbf{Latex}: The report is formatted using LaTeX for professional presentation.
\end{itemize}

% Training Section
\section{Training}
The training process involved the following steps for both the logistic regression and random forest models:

\begin{itemize}
    \item \textbf{Data Loading}: The JSON dataset was loaded using pandas and processed to extract features and labels.
    \item \textbf{Data Splitting}: The dataset was split into:
        \begin{itemize}
            \item Training set: 70\% of the data, including 836 samples.
            \item Validation set: 15\% of the data, including 179 samples.
            \item Test set: 15\% of the data, including 180 samples.
        \end{itemize}
        Splitting was performed using Scikit-learn's \texttt{train\_test\_split} with a random state of 42 for reproducibility.
    \item \textbf{Training Approach}:
                \begin{itemize}
                    \item Loss Function: Cross-Entropy Loss
                    \item Optimizer: Adam with learning rate of 0.003
                    \item Weight Decay: 1e-5 for L2 regularization
                    \item Early Stopping: Implementation of patience (16 epochs) to prevent overfitting
                    \item Pocket Algorithm: Saves the best model based on validation accuracy
                    \item Maximum Epochs: 64
                \end{itemize}
\end{itemize}

% Evaluation Metrics Section
\section{Evaluation Metrics}
The model’s performance was evaluated using the following metrics:

\begin{itemize}
    \item \textbf{Accuracy}: The proportion of correct predictions on the validation and test sets.
\end{itemize}

% Accuracy Report Section
\section{Accuracy Report}
The performance of our models is reported below:

\begin{itemize}
    \item \textbf{Phase 1: Traditional Machine Learning Models}:
        \begin{itemize}
            \item \textbf{Logistic Regression}:
                \begin{itemize}
                    \item \textbf{Validation Accuracy}: 0.9497
                    \item \textbf{Test Accuracy}: 0.9167
                    \item \textbf{Kaggle Score}: 0.9334
                \end{itemize}
            \item \textbf{Random Forest}:
                \begin{itemize}
                    \item \textbf{Validation Accuracy}: 0.8939
                    \item \textbf{Test Accuracy}: 0.9000
                \end{itemize}
        \end{itemize}
    \item \textbf{Phase 2: Neural Network Approach}:
        \begin{itemize}
            \item \textbf{Multilayer Perceptron (MLP)}:
                \begin{itemize}
                    \item \textbf{Validation Accuracy}: 0.9385
                    \item \textbf{Test Accuracy}: 0.9611
                    \item \textbf{Kaggle Score}: 0.9411
                \end{itemize}
        \end{itemize}
\end{itemize}

The MLP model achieved its best validation accuracy of 93.85\% after just 3 epochs of training, with early stopping triggered after 16 epochs without improvement. The final test accuracy was 96.11\%, indicating a significant improvement over the traditional models.

% Submission File Section
\section{Submission}
Submission files include:
\begin{itemize}
    \item \texttt{Group30\_Turkey30.ipynb}: The Jupyter notebook containing the code for data preprocessing, model training, and evaluation.
    \item \texttt{This report}: The LaTeX report detailing the project, methodology, and results.
    \item \texttt{submission.csv}: Our neural network approach achieved a higher score of 94.11\% on Kaggle.
\end{itemize}

% Conclusion Section
\section{Conclusion}

\indent In this project, we successfully implemented a machine learning model using MLP in Pytorch framework. Although the Kaggle score result did not improve significantly compared to the logistic regression model, we gained valuable insights into the classification of audio embeddings. The project demonstrated the effectiveness of traditional machine learning models and neural networks in handling audio data.

% References Section
\section{References}
\begin{itemize}
    \item Pytorch tutorial: \url{https://docs.pytorch.org/tutorials/beginner/blitz/cifar10_tutorial.html}
    \item Pandas Documentation: \url{https://pandas.pydata.org/}
    \item Google Colab: \url{https://colab.research.google.com/}
    \item Kaggle competition page: \\
        \url{https://www.kaggle.com/competitions/introduction-to-machine-learning-project-cq-24-2}
\end{itemize}

\section{AI Usage}
The project utilized AI tools including GPT-4, Grok, and Claude Sonnet 3.7 for various tasks such as:

\begin{itemize}
    \item Code syntax support and debugging.
    \item Report helping and formatting in LaTeX assistance.
\end{itemize}

We acknowledge the use of AI tools in enhancing our productivity and code quality, while ensuring that the core logic and implementation were developed by the team.\\
\indent Chat conversations with AI can be found here:
\begin{itemize}
    \item \url{https://grok.com/share/bGVnYWN5_b1b4cc44-6408-4ae0-ad5e-ed776476da01}
\end{itemize}

\end{document}